\section{Covers, isogenies and other transfers}\label{sec:transfer}

A recurring hope was that some other representation of the same
instance---an isogenous curve, a curve at another level of the isogeny
volcano, a Jacobian, a Weil descent---would make index calculus cheaper.
These routes produced the programme's clearest negatives, and one lemma
that explains most of them.

\begin{lemma}[transport conservation,
\art{docs/ic/RESEARCH_TRANSFER_AWARE_IC_20261010.md} \S1]\label{lem:conservation}
Let $\varphi:E\to E'$ be an $\FF_q$-rational isogeny (or any homomorphism
of rational points), $\kappa=\ker\varphi\cap E(\FF_q)$, $R\in E(\FF_q)$,
$\Fb'\subset\varphi(E(\FF_q))$ a factor base on the codomain and
$\Fb=\varphi^{-1}(\Fb')\cap E(\FF_q)$. Then the number of ordered
$m$-decompositions of $R$ over $\Fb$ is $\abs\kappa^{m-1}$ times the number
of ordered $m$-decompositions of $\varphi(R)$ over $\Fb'$.
\end{lemma}
\begin{proof}
Images of decompositions of $R$ are decompositions of $\varphi(R)$.
Conversely, for $\sum Q_i=\varphi(R)$ choose rational preimages $P_i$ of
the $Q_i$; then $\sum P_i-R\in\kappa$, and of the $\abs\kappa^m$ choices
exactly $\abs\kappa^{m-1}$ sum to $R$ (translate $P_1$ by the kernel
element).
\end{proof}

So a homomorphic transfer changes nothing but three curve-specific inputs:
the rational kernel symmetry $\abs\kappa^{m-1}$ (bounded by the rational
torsion; trivial on prime-order curves), the factor-base geometry
available on the destination, and the cost of arithmetic there. Sequences
of isogenies compose to one homomorphism; a loop is an endomorphism and
acts on the subgroup as a known scalar, useful only as a known-log
multiplier. Multi-curve factor bases glued by isogenies are single-curve
bases in disguise.

\paragraph{Conductor gaps in the ECC2K-130 class.} Jao, Miller and
Venkatesan~\cite{jmv2005} prove random self-reducibility among curves with
the same endomorphism ring, leaving open whether curves whose conductors
differ by a large prime could differ in hardness. The ECC2K-130 class has
conductor gap 263; we computed its structure
(\art{docs/ic/RESEARCH_ISOGENY_CONDUCTOR_GAP_20261007.md}) and ran
several rounds of experiments. Native counted-$S$ index calculus on two fixed CM
pairs gave floor/crater ratios $0.975$ (99\% $[0.603,1.451]$) at
$\ell=8209$ and $1.127$ $[0.922,1.342]$ at $\ell=131101$: no level effect
(192 verified records, 96 exact replays;
\art{research/ecbench_conductor_gap_20261012/}). Two-torsion
symmetrisation is a quotient isogeny: at a matched base of abscissae it
cut F4 row-reduction work $4.1$--$4.3\times$ ($\abs H=2$) and
$17.6$--$19.6\times$ ($\abs H=4$) over 1{,}128 systems---but a correction
the next day showed that at a matched number of \emph{columns} the saving
disappears; it was the cost of a redundant orbit base. A combined Weil
descent used $234$--$529\times$ the F4 operations of the blind system
(\art{research/isogeny_conductor_gap_round3_20261011/}).

\paragraph{Fibre-aware index calculus.} Isogeny fibres are torsors whose
arithmetic is controlled by one cohomology class
(\art{RESEARCH_ISOGENY_FIBERS_ECDLP.md}); relation generation is
isogeny-equivariant at zero lifting cost, and a fibre-closed pair table was
the best lever of the conductor-gap rounds ($12.8\times$ rho in its best
cell). Measured against matched rho across sizes at fixed cofactor, the
envelope IC/rho ratio grows from $8.4$--$8.8$ to $35.1$--$37.0$ ($h=4$) and
from $13.4$--$20.1$ to $60.4$--$63.1$ ($h=72$), as $\#E^{1/6}$: the gap
widens, it does not close (816 audited records;
\art{research/fiber_ic_scaling_20261011/}).

\paragraph{Other structure.} Symmetrised $u$-frame bases saved
$0.84\times$ at $n=17$ and $0.76$--$0.77\times$ at $n=23$, but cost
$4.01\times$ at $n=31$ (\art{research/koblitz_symmetrised_e2e_20260927/}).
GHS and hyperelliptic covers: ECC2K-130 does map to hyperelliptic curves,
but the genus that would pay is missing. In odd characteristic the cubic
GHS transfer from $E_0(\FF_{q^3})$ to a genus-3 Jacobian now runs end to end
with double-large-prime index calculus, at about $q^{1.4}$ against rho's
$q^{3/2}$, within $1.15\times$ of a measured negation-map rho at
$q=2^{20}$ (\S\ref{sec:ladder}). The
isogeny-loop GLV/GLS work produced a measured $1.09$--$1.15\times$ faster
four-dimensional scalar multiplication (\art{research/isogeny_loop_glv_20261012/})
and a $\sqrt{\text{\'elu}}$ kernel cost exponent $0.822$ $[0.805,0.838]$;
these speed up curve arithmetic, not the discrete logarithm.

\paragraph{Covers, 21--24 October.} The Joux--Vitse genus-3 cover of
$E(\FF_{p^6})$ is now within $1.6$--$2.0\times$ of a pooled rho constant at
$p=1009$; measuring the per-trial cost on a genus-255 binary Jacobian
turned the $N=130$ GHS counterfactual's modelled margin negative ($-1.83$
bits); and two censuses bound what covers exist at toy sizes (genus floor 13
for odd-order $E/\FF_{q^4}$; a $2^{-l}$ genus-3 fraction over
$\FF_{2^{3l}}$). All are in \S\ref{sec:ladder3}. None touches a prime field
or a prime-degree binary field, where the GHS genus is out of reach.

\paragraph{Genus three against the strongest rho.} Rounds three and four
of the genus-3 work ran without a rho arm; a 23 October study supplied one
(\art{research/ghs3_vs_rho_20261023/TABLES.md}; the directory has a
protocol and tables but no README). On round 4's instances, $q=2^{14}$ to
$2^{20}$, three curves per size, in one unit and one binary, the
preregistered per-size decision is \emph{IC below the strong rho} at
$q=2^{16}$ (geometric mean IC/$R_{\mathrm{ideal}}$ $0.701$) and
\emph{unresolved} at the other six sizes (geometric means $0.97$--$1.52$;
$1.274$ at $2^{20}$). The ratio carries $\sqrt c$ of the cofactor, because IC
depends on $q$ alone and rho on $\ell=\#E_0/c$; a post-hoc fit attributes the
sub-rho cells to cofactor 4. Pricing the subgroup multiplication at
$c_\ell=9$ instead puts every instance at $2.7$--$12.4\times$ rho, and the fitted
trend is flat ($0.034$ $[-0.047,0.105]$ per bit), so the extrapolated
crossover is undetermined ($\log_2q=11.5$, interval $[-35.6,60.4]$). Genus
three after cubic GHS is therefore at parity with rho to within the pricing
of one operation on these sizes, not past it. The class is structured
($E_0/\FF_{q^3}$ with $4\mid\#E_0$), not prime-order.

\paragraph{Joux--Vitse: two generations of comparison.} The current
reading is the 24 October $p$-scaling study, $1.59$--$1.98\times$ a pooled rho
constant at $p=1009$ and closing as about $p^{-2}$ (\S\ref{sec:ladder3}).
Earlier readings were lower and are not current. A sieve variant measured
$0.024$ $[0.019,0.031]$ of a strong rho at $p=1009$ and $0.0042$ at
$p=1823$, but its 20 October review withheld every comparison (no numeric
acceptance for the controls, uncharged rho hashes, incomplete accounting;
\art{research/jv_sieve_strong_rho_20261019/}); a prospective re-run is
registered and has not run. A 9 October whole-route replay from a non-weak
challenge curve, through an isogeny walk to a weak curve and its cover, read
$0.058$ of a pooled rho at $p=1009$ and $0.011$ at $p=1511$, on one seed per
prime, with the historical method revision and an unmatched rho
(\art{research/jv_cover_e2e_30bit_20261009/RESULT.md}). The cover-decomposition
note still states a measured sieve crossover near $p\approx430$
(\art{research/notes/index-calculus/RESEARCH_COVER_DECOMPOSITION_LEDGER.md}
\S11) and carries no erratum; it should be read through the review.
Leaderboard Table~D, the historical JV study, is withheld. On the method
side, a sparse-FGLM linear stage made the cover $1.28\times$ cheaper end to
end (\art{research/jv_cover_sparse_fglm_20261018/}); relation collection is
$99.99\%$ of $S$ (\art{research/ic_cover_leads_20261018/}); depth-1 isogeny
classes hold no weak curve over 4{,}871{,}062 curves, while 86--95\% of
deeper classes hold one (\art{research/jv_weak_class_law_20261007/}); and
the $(n-1)$-decomposition variant is $4.5\cdot10^7$--$1.8\cdot10^8\times$
rho (\art{research/joux_vitse_n1_20261015/}).

\paragraph{Cover censuses.} The minimal cubic GHS genus over an isogeny
class follows $\#E(K)\bmod4$: 13 for odd order, 5 for $2\bmod4$, 3 for
$0\bmod4$, and genus 3 is the only value below rho's $q^{3/2}$
(\art{research/ghs_cubic_cover_census_20261012/}); at $n=5$ the floors are
33, 33, 25 and 5 by 2-torsion type (\art{research/ghs_cover_census_n5_20261013/}).
No deployed prime-field or prime-degree binary curve has a usable subfield;
only the RFC~2409 Oakley groups 3 and 4 are reachable, at genus 15--16, where
exponent models disagree in sign and the measured genus-16 relation phase
alone is 5.2 bits slower than rho (\art{research/cover_route_prime_binary_20261008/},
\art{research/ghs_oakley_extended_20261023/}). All 15{,}791 cover maps in the
registry were assessed: 165 binary maps have genus-3 arithmetic and the
15{,}626 odd-characteristic genus-2 covers lack a Jacobian implementation
(\art{research/cover_transfer_assessment_20261013/}). An inventory of about
90 cover and decomposition studies found every route on prime, Koblitz and
prime-degree binary curves closed, negative or constant-factor
(\art{research/cover_decomposition_gap_search_20261018/}). Smith's
$(2,2,2)$-isogeny route was reproduced natively (full success $0.18638$
against a theoretical $0.18574$ over $10^6$ curves) and its sign bias shrinks
like $1/\sqrt p$ (\art{research/smith_prop24_native_20261014/},
\art{research/smith_sign_scaling_20261016/}).

\paragraph{Volcano levels and generators.} A 403{,}168-run census found the
level effect equivalent in 33 of 35 comparisons and different in none
(\art{research/volcano_ecdlp_census_20261017/}), and a large-sample rho study
found the floor/crater ratio $0.998$ $[0.986,1.010]$
(\art{research/conductor_gap_rho_large_sample_20261018/}). The one real level
effect is the automorphism group: $\sqrt3$ on the $j=0$ crater, confirmed at
$1.86$--$1.92\times$ (\art{research/drastic_prime_conductor_gap_round3_20261012/}).
At $\ell=2$, floors give $5.8\%$ fewer relations than craters, and 99\% of the
contrast is explained by the base's cofactor classes
(\art{research/isogeny-volcano-cost-v2_20261017/},
\art{research/generator_effect_ic_20261019/}).

\paragraph{Named obstructions.} Blocked requirements are kept open as
problem records (\art{problems/README.md}), among them the complete-cost
calibration of root-finding IC
(\art{problems/rootfinding-complete-cost-calibration.md}), with exit
criteria rather than abandoned.
